\documentclass[11pt]{article}
\usepackage[utf8]{inputenc}
\usepackage[T1]{fontenc}
\usepackage{lmodern}
\usepackage[a4paper,margin=2.2cm]{geometry}
\usepackage{amsmath,amssymb}
\usepackage{booktabs}
\usepackage{array}
\usepackage{xcolor}
\usepackage{titlesec}
\usepackage{enumitem}
\usepackage[hidelinks]{hyperref}
\definecolor{navy}{HTML}{1A1A1A}
\titleformat{\section}{\normalfont\large\bfseries\color{navy}}{\thesection.}{0.5em}{}
\titleformat{\subsection}{\normalfont\normalsize\bfseries\color{navy}}{\thesubsection.}{0.5em}{}
\titlespacing{\section}{0pt}{8pt}{4pt}
\titlespacing{\subsection}{0pt}{6pt}{2pt}
\setlist{nosep,leftmargin=1.4em}
\newcommand{\USD}{\mathrm{USD}}
\pagestyle{plain}

\begin{document}
\begin{center}
{\Large\bfseries\color{navy} Impacto económico del RIGI: nota metodológica}\\[2pt]
{\large Modelo insumo--producto provincial para la estimación del impacto del primer año}\\[4pt]
{\small Estimación de producto, valor agregado y empleo (formal e informal) por provincia}\\[2pt]
{\normalsize\bfseries Pedro Elosegui}\\[2pt]
{\small\itshape Septiembre 2026}
\end{center}
\vspace{-4pt}
\hrule
\vspace{6pt}

\section{Objetivo y fuentes}
Se estima el impacto de corto plazo (primer año) de los proyectos adheridos al Régimen de Incentivo a las Grandes Inversiones (RIGI) sobre el producto, el valor agregado (VA) y el empleo, con apertura por provincia y por condición de formalidad. El ejercicio es de \emph{demanda}, utilizando una matriz insumo-producto por provincias (MIPR), y toma como shock de demanda la inversión efectivamente comprometida, no la anunciada.

Las fuentes son: (i) \textbf{Monitor RIGI-IIEP} (\texttt{base\_completa.xlsx}), del que se toma la \emph{inversión comprometida en activos computables del Año 1} por proyecto; (ii) \textbf{MIP Provincial Argentina 2023} (Michelena, G., Herrera Gómez, M.\ y Elosegui, P.; $24$ provincias $\times\,20$ sectores $=480$ actividades), con vectores de VBP, VAB y empleo; (iii) \textbf{ENGHo 2017--18} para el patrón de consumo de los hogares; (iv) \textbf{matriz de comercio interprovincial} para la localización del efecto inducido; y (v) tasas de \textbf{formalidad} por provincia y sector (empleo registrado sobre total).

\section{El shock: función de producción por proyecto}
La inversión inicial no se imputa a un único sector. Para cada proyecto $k$, con inversión de primer año $I_k$ (en USD) localizada en el conjunto de provincias $\mathcal{P}_k$ y categoría tecnológica $c(k)$, el gasto se descompone en cuatro rubros con participaciones $(\alpha_c,\beta_c,\gamma_c,\delta_c)$, $\alpha_c+\beta_c+\gamma_c+\delta_c=1$:
\[
\underbrace{\alpha_c}_{\text{obra civil}}\!\to\text{Construcción},\quad
\underbrace{\beta_c}_{\text{equipo nac.}}\!\to\text{Máq. y equipos},\quad
\underbrace{\gamma_c}_{\text{servicios}}\!\to\text{Serv. empresariales},\quad
\underbrace{\delta_c}_{\text{importado}}\!\to\text{fuga}.
\]
La porción importada $\delta_c$ no genera impacto doméstico. El vector de demanda final (en pesos del año base) reparte cada proyecto en partes iguales entre sus provincias:
\begin{equation}
\mathbf{f} \;=\; \tau \sum_{k}\; \frac{I_k}{|\mathcal{P}_k|}\sum_{p\in\mathcal{P}_k}\Big(\alpha_{c(k)}\,\mathbf{e}_{p,\text{con}}+\beta_{c(k)}\,\mathbf{e}_{p,\text{maq}}+\gamma_{c(k)}\,\mathbf{e}_{p,\text{serv}}\Big),
\end{equation}
donde $\mathbf{e}_{p,s}$ es el vector unitario que selecciona la celda (provincia $p$, sector $s$) y $\tau$ el tipo de cambio del año base (véase §7).

\section{Propagación: producto y valor agregado}
Con la matriz de coeficientes técnicos $\mathbf{A}$ (columna $j$: $a_{ij}=z_{ij}/x_j$), la inversa de Leontief $\mathbf{L}=(\mathbf{I}-\mathbf{A})^{-1}$ propaga el shock:
\begin{equation}
\mathbf{x}=\mathbf{L}\,\mathbf{f},\qquad
\mathbf{x}^{\text{dir}}=\mathbf{f},\qquad
\mathbf{x}^{\text{ind}}=\mathbf{x}-\mathbf{f},\qquad
\mathbf{va}=\mathbf{v}\odot\mathbf{x},
\end{equation}
con $\mathbf{v}$ el vector de coeficientes VAB/VBP por celda. El radio espectral de $\mathbf{A}$ es $\rho(\mathbf{A})\approx0{,}44$, lo que garantiza la convergencia de $\mathbf{L}=\sum_t \mathbf{A}^t$. Se supone oferta perfectamente elástica y coeficientes técnicos fijos (sin cuellos de botella ni desplazamiento de otra producción).

\section{Efecto inducido: cierre Tipo II por la orla de hogares}
El modelo se cierra incorporando el sector hogares a la matriz, con una \emph{orla} de $24$ hogares (uno por provincia). La matriz ampliada $\mathbf{A}^{\ast}$ agrega, a los $N=480$ sectores, una fila y una columna por hogar:
\begin{itemize}
\item \textbf{Fila (generación del ingreso):} salario generado por unidad de producción de cada celda de la provincia $p$, $\;a^{\ast}_{h_p,\,ps}=\omega_{ps}\,v_{ps}$, con $\omega_{ps}=\mathrm{RTA}_{ps}/\mathrm{VAB}_{ps}$ la participación salarial de la celda.
\item \textbf{Columna (consumo):} gasto por unidad de ingreso del hogar $p$, $\;a^{\ast}_{ps,\,h_p}=\pi\,s_s\,T_{p,\,\mathrm{prov}(ps)}$, con $\pi=0{,}75$ la propensión marginal a consumir, $s_s$ el patrón sectorial de consumo (ENGHo, excluye administración pública) y $\mathbf{T}$ la matriz de comercio interprovincial fila-normalizada que rutea el gasto a la provincia de producción.
\end{itemize}
La inversa ampliada $\mathbf{L}^{\ast}=(\mathbf{I}-\mathbf{A}^{\ast})^{-1}$ genera el efecto inducido de forma \emph{endógena}: producción, valor agregado y empleo surgen del mismo circuito ingreso$\to$consumo$\to$producción, sin coeficientes pegados. El radio espectral $\rho(\mathbf{A}^{\ast})=0{,}66<1$ garantiza convergencia.

\paragraph{Imputación de la participación salarial.} La MIP provincial trae una sola fila de VAB; su apertura en salario (RTA) y excedente$+$ingreso mixto se imputa como $\omega_{ps}=\min\!\big(\omega_s\,a_{ps}/\bar a_s,\;c\big)$, con $c=0{,}85$, donde $\omega_s$ es la participación salarial nacional del sector (Cuenta de Generación del Ingreso) y $a_{ps}$ la participación de asalariados (registrados $+$ no registrados) en el empleo de la celda (de \texttt{empleo\_prov\_sec\_ipf}). Se renormaliza dentro de cada sector para que $\sum_p \omega_{ps}\mathrm{VAB}_{ps}=\omega_s\sum_p\mathrm{VAB}_{ps}$ (coherencia con Cuentas Nacionales). Así $\omega_{ps}\in[0;\,0{,}85]$, mayor donde el sector es más asalariado.

\section{Empleo: nuevos puestos del primer año (marginal)}
Aplicar los coeficientes de empleo a toda la producción ($\mathbf{e}^{\top}\mathbf{L}^{\ast}\mathbf{f}$) mide el empleo \emph{embebido} —cuánta gente sostiene esa producción a intensidad media—, un concepto de estado estacionario que sobreestima los puestos \emph{nuevos} de un pulso de obra de un año (y $\sim\!15\times$ en el directo, porque la obra de un megaproyecto es mucho más intensiva en capital que la construcción promedio de la matriz; véase Anexo B). Se estima entonces la \emph{movilización marginal}:
\begin{align}
E^{\text{dir}} &= \sum_k \tfrac{I_k}{|\mathcal{P}_k|}\big(\lambda^{o}_{c(k)}+\lambda^{p}_{c(k)}\big) && \text{(directo, por intensidades de ingeniería)}\\
E^{\text{ind}} &= (\mu_{\mathrm{I}}-1)\,E^{\text{dir}}, & \mu_{\mathrm{I}}&=\mathbf{e}^{\top}\mathbf{L}\mathbf{f}\,/\,\mathbf{e}^{\top}\mathbf{f}=1{,}38\\
E^{\text{indu}} &= (\mu_{\mathrm{II}}-\mu_{\mathrm{I}})\,E^{\text{dir}}, & \mu_{\mathrm{II}}&=\mathbf{e}^{\top}\mathbf{L}^{\ast}\mathbf{f}\,/\,\mathbf{e}^{\top}\mathbf{f}=1{,}98
\end{align}
Los multiplicadores de empleo $\mu_{\mathrm{I}},\mu_{\mathrm{II}}$ \emph{salen del modelo} cerrado, no se fijan a mano. El indirecto y el inducido se distribuyen entre celdas con los coeficientes de empleo aplicados a la producción de cada ronda ($\mathbf{x}^{\text{ind}},\mathbf{x}^{\text{indu}}$); el total es la suma de las celdas provincia-sector. El directo de obra se asigna a Construcción y el operativo al sector propio del proyecto.

\paragraph{Formalidad.} Sea $\boldsymbol{\phi}$ la tasa de formalidad por celda (registrado/total). La obra directa se toma formal; el resto se pondera por $\boldsymbol{\phi}$:
\begin{equation}
E^{\text{formal}} = E^{\text{dir,obra}} + \boldsymbol{\phi}\odot\big(E^{\text{dir,op}}+E^{\text{ind}}+E^{\text{indu}}\big),\qquad
E^{\text{informal}}=E-E^{\text{formal}}.
\end{equation}

\section{Neutralidad del tipo de cambio}
El tipo de cambio del año base $\tau$ (296 ARS/USD, 2023) sólo traduce el shock a la unidad de la matriz. Como el VA se reporta en dólares dividiendo por el \emph{mismo} $\tau$, éste se cancela:
\begin{equation}
\mathrm{VA}^{\USD}=\frac{\mathbf{v}^{\!\top}\mathbf{L}\,\mathbf{f}}{\tau}=\frac{\mathbf{v}^{\!\top}\mathbf{L}\,(\tau\,\mathbf{I}^{\USD})}{\tau}=\mathbf{v}^{\!\top}\mathbf{L}\,\mathbf{I}^{\USD}.
\end{equation}
El empleo, anclado en dólares, también es invariante a $\tau$. El supuesto sustantivo no es el nivel del dólar sino la \textbf{estructura productiva y de precios relativos de 2023}, que se asume constante hacia 2026.

\section{Resultados de referencia (primer año)}
Con la inversión comprometida de Año 1 de los 22 proyectos aprobados ($5.511$ millones de USD):

\begin{center}\small
\begin{tabular}{lr}
\toprule
\textbf{Agregado (primer año)} & \textbf{Valor} \\
\midrule
Producto (USD millones) & 9.470 \\
Valor agregado (USD millones) & 4.796 \\
Empleo total (puestos) & 12.010 \\
\quad directo & 6.055 \\
\quad indirecto & 2.329 \\
\quad inducido & 3.627 \\
Empleo formal (puestos, 69\%) & 8.312 \\
Empleo informal (puestos, 31\%) & 3.698 \\
\bottomrule
\end{tabular}
\end{center}

\paragraph{Banda de sensibilidad al multiplicador.} El empleo escala con el multiplicador Tipo II $\mu_{\mathrm{II}}$; se reporta la banda para $\mu_{\mathrm{II}}\in[1{,}80;\,2{,}18]$ (central $1{,}98$):
\begin{center}\small
\begin{tabular}{lccc}
\toprule
$\mu_{\mathrm{II}}$ & 1{,}80 & 1{,}98 (central) & 2{,}18 \\
\midrule
Empleo total (puestos) & 10.900 & 12.010 & 13.200 \\
\bottomrule
\end{tabular}
\end{center}
El producto y el VA no dependen de $\mu_{\mathrm{II}}$.

%\paragraph{Distribución provincial (principales).} Río Negro concentra el mayor impacto (2.505 puestos; 84\% formal; 0{,}80\% del empleo provincial), seguida por Buenos Aires (2.020; 61\%; 0{,}03\%), San Juan (1.761; 77\%), Neuquén (1.332; 77\%) y Salta (1.152; 74\%). El detalle completo por las 24 jurisdicciones se acompaña en planilla.

\appendix
\section{Anexo A --- Parámetros de composición por tipo de proyecto}
Participaciones del CAPEX $(\alpha,\beta,\gamma,\delta)$ e intensidades de empleo $(\lambda^o,\lambda^p)$ y formalidad operativa por categoría:
\begin{center}\footnotesize
\begin{tabular}{lcccccc c}
\toprule
Categoría & Obra $\alpha$ & Eq.\,nac.\,$\beta$ & Import.\,$\delta$ & Serv.\,$\gamma$ & $\lambda^o$ & $\lambda^p$ & Form.\,op. \\
\midrule
Transporte de gas (PG\_TRANS) & 0{,}52 & 0{,}13 & 0{,}25 & 0{,}10 & 1{,}0 & 0{,}2 & 0{,}94\\
GNL & 0{,}15 & 0{,}05 & 0{,}72 & 0{,}08 & 0{,}15 & 0{,}1 & 0{,}94\\
\emph{Upstream} (PG\_UP) & 0{,}28 & 0{,}12 & 0{,}38 & 0{,}22 & 0{,}3 & 0{,}2 & 0{,}94\\
Minería metalífera & 0{,}30 & 0{,}05 & 0{,}29 & 0{,}36 & 0{,}8 & 0{,}3 & 0{,}66\\
Minería (lixiviación) & 0{,}42 & 0{,}05 & 0{,}20 & 0{,}33 & 0{,}8 & 0{,}3 & 0{,}66\\
Litio & 0{,}50 & 0{,}08 & 0{,}22 & 0{,}20 & 0{,}8 & 0{,}5 & 0{,}66\\
Energía eléctrica & 0{,}30 & 0{,}10 & 0{,}55 & 0{,}05 & 0{,}8 & 0{,}2 & 0{,}74\\
Siderurgia & 0{,}30 & 0{,}12 & 0{,}48 & 0{,}10 & 1{,}0 & 0{,}5 & 0{,}58\\
Infraestructura & 0{,}62 & 0{,}08 & 0{,}20 & 0{,}10 & 1{,}2 & 0{,}3 & 0{,}50\\
\bottomrule
\end{tabular}
\end{center}
\noindent Las participaciones se calibran con estudios de factibilidad (NI 43-101) y estructura típica de CAPEX por tecnología. La fuga importada oscila entre 20\% (infraestructura, litio) y 72\% (GNL).

\section{Anexo B --- Por qué no se usa el empleo bruto de la matriz}
Aplicar los coeficientes de empleo a la producción propagada arroja, para este shock, \textbf{131.079 puestos} (directo+indirecto), frente a los \textbf{8.446} del empleo anclado: un factor de \textbf{15{,}5$\times$}. El $57\%$ de ese empleo bruto proviene de un único sector, Construcción, por tres razones:
\begin{enumerate}
\item \textbf{Promedio aplicado a lo marginal.} El coeficiente es la intensidad \emph{media}: el trabajador promedio de construcción en la matriz produce $4,77$ millones de pesos de VA al año ($\approx \mathrm{USD}~16.000$), típico de obra informal de baja productividad. La obra de un tren de GNL, un mineraloducto o el desarrollo de una mina la ejecutan empresas especializadas, intensivas en capital y de alta productividad; la intensidad media no aplica.
\item \textbf{Dispersión de productividad.} El VA por puesto varía entre $1,96$ y $469$ millones de pesos según el sector ($\sim\!240\times$). El resultado se vuelve hipersensible a hacia dónde $\mathbf{L}$ rutea el producto, y lo rutea hacia los sectores más intensivos en trabajo (construcción, comercio, servicios personales).
\item \textbf{Stock vs. flujo y regionalización.} El coeficiente refleja la dotación total que sostiene un nivel de producción (incluye trabajo fijo, \emph{overhead} e informal), no la contratación marginal; además, la matriz provincial se construye regionalizando la tabla nacional con cocientes de localización, mientras el empleo proviene de otra fuente, lo que introduce desajustes que se amplifican al invertir.
\end{enumerate}
\noindent Por ello $\mathbf{e}^{\top}\mathbf{L}\mathbf{f}$ se interpreta como \emph{cota superior} y el empleo directo se ancla a parámetros de ingeniería, con multiplicadores acotados para el indirecto e inducido.

\vspace{4pt}
\hrule
\vspace{3pt}
{\footnotesize\itshape Pedro Elosegui (IIEP-UBA), sobre Monitor RIGI--IIEP y MIP provincial 2023. Modelo de demanda (insumo--producto), equilibrio parcial: no incorpora precios relativos endógenos, restricciones de recursos, desplazamiento (\emph{crowding-out}) ni impacto fiscal o de divisas. Estima el primer año (fase de obra), no la producción futura de los yacimientos.}

\end{document}
